\chapter{Jacobi fields on a Delaunay unduloid}\label{sec:delaunay-jacobi}
The geometric Jacobi fields and their growth are described in \cite{Maz02}, Section 2.3, pp.~8--10. For the noncylindrical unduloids of Theorem \ref{thm:explicit-unduloid-construction}, we prove a bounded-kernel statement by explicit calculations and a one-dimensional barrier. Weighted Fredholm theory and precise Floquet exponents remain outside this section.

\begin{proposition}[The axial period is strictly increasing with neck radius]\label{prop:unduloid-period-derivative}
For fixed $H>0$ and $0<a<1/H$, the axial period $L_a$ in Proposition \ref{prop:unduloid-global-geometry} has $L_a'>0$. In particular $4/H<L_a<2\pi/H$.
\end{proposition}
\begin{proof}
Set $e=1-Ha\in(0,1)$ and $t=Hs$. The radius formula and $2Hc=1-e^2$ give
\[
 L_a=\frac1H\int_0^{2\pi}\frac{1-e\cos t}{\sqrt{1+e^2-2e\cos t}}\,dt.
\]
For a fixed $e<1$ the denominator is uniformly positive, so differentiation under the integral is valid. Direct differentiation of the integrand with respect to $e$ gives
\[
 -\frac{e\sin^2t}{(1+e^2-2e\cos t)^{3/2}}.
\]
Since $de/da=-H$, it follows that
\[
 L_a'=\int_0^{2\pi}\frac{e\sin^2t}{(1+e^2-2e\cos t)^{3/2}}\,dt>0.
\]
Combine strict monotonicity with the two endpoint limits in Propositions \ref{prop:unduloid-cylinder-limit} and \ref{prop:unduloid-sphere-chain-limit} to get the bounds.
\end{proof}

\begin{proposition}[Six geometric Jacobi fields]\label{prop:unduloid-geometric-jacobi}
Write $r=r_a$, $z=z_a$, let primes denote arclength derivatives, and let $\partial_a$ differentiate the fixed-$H$ family at fixed $s$. The following six functions are Jacobi fields for $J=\Delta+|A|^2$ on a noncylindrical unduloid:
\[
 -r',\quad z'\cos\theta,\quad z'\sin\theta,\quad
 (zz'+rr')\cos\theta,\quad(zz'+rr')\sin\theta,\quad
 f_a=(\partial_ar)z'-(\partial_az)r'.
\]
They are linearly independent. The first three come from translations, the next two from rotations of the axis, and the last from changing the neck radius at fixed $H$.
\end{proposition}
\begin{proof}
For any smooth family of immersions with the same constant $H$, the normal speed satisfies $Jf=0$ by Proposition \ref{prop:riemannian-mean-curvature-variation}; tangential speed differentiates a constant and contributes zero. Pair the three constant ambient vector fields with $\nu=(z'\cos\theta,z'\sin\theta,-r')$ to get the translation fields. Rotation about the $y$ axis has velocity $(z,0,-r\cos\theta)$ and hence the fourth normal component; rotation about the $x$ axis gives the fifth up to sign. Pair $\partial_aF$ with $\nu$ to get the sixth.

The angular zero modes $-r'$ and $f_a$ are independent: at the neck $s=0$, they take values zero and one respectively, while $r'$ is not identically zero. For each first angular mode, the translation coefficient $z'$ is positive and periodic. The rotational coefficient grows without bound because $z$ grows linearly, $z'$ is bounded below positively, and $rr'$ is bounded. They are therefore independent. Fourier orthogonality separates the zero, cosine and sine sectors, proving independence of all six fields. Rotation around the symmetry axis has identically zero normal speed and is not counted.
\end{proof}

\begin{proposition}[The neck Jacobi field has linear drift]\label{prop:unduloid-neck-jacobi-drift}
With $P=2\pi/H$ and $f_a$ as above,
\[
 f_a(s+P)=f_a(s)-L_a'r'(s).
\]
Consequently $f_a$ is unbounded on the complete noncylindrical unduloid, and no combination of $f_a$ with the vertical translation field can be bounded unless the coefficient of $f_a$ vanishes.
\end{proposition}
\begin{proof}
Differentiate $r_a(s+P)=r_a(s)$ and $z_a(s+P)=z_a(s)+L_a$ with respect to $a$. The period $P$ is independent of $a$, so no derivative of the argument is introduced. Substitution into the formula for $f_a$ proves the drift relation. At a value of $s$ where $r'(s)\ne0$, iteration gives $f_a(s+jP)=f_a(s)-jL_a'r'(s)$. The coefficient is nonzero by Proposition \ref{prop:unduloid-period-derivative}, so this grows linearly. The vertical translation field is periodic and cannot cancel this drift.
\end{proof}

\begin{lemma}[A bounded solution of a strictly positive one-dimensional potential vanishes]\label{lem:positive-ode-no-bounded-solution}
Let $q:\mathbb R\to\mathbb R$ be continuous with $q\ge\delta>0$. If a bounded real $C^2$ function $v$ satisfies $v''=qv$ on $\mathbb R$, then $v=0$.
\end{lemma}
\begin{proof}
Let $M=\sup|v|$. On $[-T,T]$ define $b_T(t)=M\cosh(\sqrt\delta\,t)/\cosh(\sqrt\delta\,T)$. For $\mathcal L=-d^2/dt^2+q$, one has $\mathcal Lb_T=(q-\delta)b_T\ge0$ and $\mathcal Lv=0$. Each of $b_T\pm v$ is nonnegative at both endpoints. If it had a negative interior minimum, its second derivative there would be nonnegative and its $q$ term strictly negative, contradicting $\mathcal L(b_T\pm v)\ge0$. Hence $|v(t)|\le b_T(t)$. For fixed $t$, let $T\to\infty$ to obtain $v(t)=0$.
\end{proof}

\begin{proposition}[The conformal Jacobi potential has a uniform bound]\label{prop:unduloid-jacobi-potential}
Let $t$ be the global coordinate defined by $dt/ds=1/r(s)$. The unduloid metric is $r^2(dt^2+d\theta^2)$, and
\[
 r^2J=\partial_t^2+\partial_\theta^2+V(t),\qquad
 V=r^2|A|^2=\frac{H^2r^2}2+\frac{2c^2}{r^2}\le2-2Hc<2.
\]
In particular every bounded Jacobi field has zero Fourier coefficients in all angular frequencies $|k|\ge2$.
\end{proposition}
\begin{proof}
Since $a\le r\le b$, the coordinate $t$ ranges over all of $\mathbb R$, and $ds=r\,dt$ gives the conformal metric. The flux equation yields principal curvatures $H/2+c/r^2$ and $H/2-c/r^2$, proving the formula for $|A|^2$. With $c=Hab/2$, the potential equals $H^2(r^2+a^2b^2/r^2)/2$. As a function of $r^2\in[a^2,b^2]$, this convex expression has its maximum at an endpoint, where it equals $H^2(a^2+b^2)/2=2-2Hc$.

A smooth bounded Jacobi field has bounded smooth Fourier coefficients $v_k(t)$, obtained by integrating against $e^{-ik\theta}$. Differentiating under this compact angular integral gives $v_k''=(k^2-V)v_k$. For $|k|\ge2$, the coefficient is at least $2+2Hc>0$. Apply Lemma \ref{lem:positive-ode-no-bounded-solution} to the real and imaginary parts to conclude $v_k=0$. No Floquet asymptotic theorem is used.
\end{proof}

\begin{theorem}[Bounded Jacobi fields on a noncylindrical unduloid]\label{thm:unduloid-bounded-jacobi}
Every smooth bounded real Jacobi field on the complete unduloid with $0<a<1/H$ is the normal component of an ambient translation. Thus the bounded kernel is exactly three-dimensional. This is the bounded part of the geometric-field description in \cite{Maz02}, p.~9; the proof below does not claim weighted Fredholm estimates or a full classification of exponentially growing fields.
\end{theorem}
\begin{proof}
Proposition \ref{prop:unduloid-jacobi-potential} eliminates all frequencies $|k|\ge2$. Fourier uniqueness leaves only the angular zero, cosine and sine sectors. In each sector the Jacobi equation is a second-order linear ODE in $t$, whose solution space has dimension two by uniqueness of the initial-value problem. Proposition \ref{prop:unduloid-geometric-jacobi} supplies two independent solutions in each sector. In the zero sector, the neck field has nonzero linear drift by Proposition \ref{prop:unduloid-neck-jacobi-drift}, while the vertical translation field is periodic. Hence boundedness leaves only the latter. In each first angular sector the rotation coefficient $zz'+rr'$ is unbounded, while $z'$ is periodic, leaving only its translation coefficient. All three retained fields are bounded and independent, completing the proof. The cylinder endpoint is excluded: axial translation has zero normal component there and additional bounded oscillatory axial fields appear.
\end{proof}

\begin{proposition}[No square-integrable Jacobi field on a complete unduloid]\label{prop:unduloid-no-l2-jacobi}
For a noncylindrical complete unduloid with fixed $0<a<1/H$, every smooth Jacobi field $u$ satisfying $\int|u|^2\,dA<\infty$ is zero. This conclusion does not assert that the Jacobi operator has closed range on an unweighted function space.
\end{proposition}
\begin{proof}
In the conformal coordinates of Proposition \ref{prop:unduloid-jacobi-potential}, $dA=r^2dt\,d\theta$ and $a\le r\le b$. Each angular Fourier coefficient $v_k$ is therefore in $L^2(\mathbb R,dt)$ by Cauchy--Schwarz on the angular circle. For $|k|\ge2$, it satisfies $v_k''=q_kv_k$ with $q_k=k^2-V\ge\delta>0$. Apply the following calculation to its real and imaginary parts. Choose $\eta_T$ equal to one on $[-T,T]$, supported in $[-2T,2T]$, with $|\eta_T'|\le C/T$. Integration by parts against $\eta_T^2v_k$ and Young's inequality give
\[
 \frac12\int\eta_T^2|v_k'|^2+\int q_k\eta_T^2|v_k|^2
 \le2\int|\eta_T'|^2|v_k|^2\le\frac{2C^2}{T^2}\|v_k\|_2^2.
\]
Letting $T\to\infty$ forces $v_k=0$ on every bounded interval.

For the remaining frequencies, the explicit bases in Proposition \ref{prop:unduloid-geometric-jacobi} apply. In a first angular sector, the rotational coefficient grows at least linearly at either end, while $z'$ is periodic and bounded below positively. No nonzero combination is square-integrable. In the zero sector, put $v=-r'$. The neck field obeys $f_a(s+jP)=f_a(s)+jL_a'v(s)$. Choose a subinterval of one period on which $|v|$ is bounded below positively. A nonzero coefficient of $f_a$ then makes the squared integral over its translates grow quadratically in $|j|$; a nonzero multiple of $v$ alone has the same positive squared integral on every period. Neither is square-integrable. The changes between $s$, $t$, and area measure have uniformly bounded positive weights for this fixed unduloid. Fourier uniqueness now gives $u=0$. This proof concerns the kernel only; range and weighted mapping properties require additional analysis.
\end{proof}
